\section{Experimental Setup}
\label{sec:setup}

\paragraph{Datasets.}
We evaluate over a 12-month rolling window spanning January--December 2024. Table~\ref{tab:dataset-matrix} summarizes the four evaluation slices. The \emph{Primary benchmark} comprises 1,847 questions drawn from fresh Reuters/BBC/AP articles and Wikipedia recent edits across all 12 months (roughly 154 questions per month). The \emph{Human-verified subset} contains 412 questions from months 3, 6, 9, and 12 that were manually adjudicated by three annotators (Cohen's $\kappa = 0.73$); this slice is used for calibration and inter-annotator agreement analysis. The \emph{Robustness slice} contains 528 questions constructed by adversarially perturbing the primary slice---replacing named entities with plausible alternatives and paraphrasing question stems---to test retrieval robustness. The \emph{Ablation slice} contains 891 questions replayed under varied system configurations for component isolation. After contamination filtering (see Section~\ref{sec:method}), we retain 1,731 primary questions (93.7\% retention rate) as clean evaluation items.

\paragraph{Retrieval and generation baselines.}
We evaluate all combinations of five retrieval systems and three generators:

\noindent\textit{Retrieval}: (1)~\textbf{BM25}~\citep{robertson2009bm25}, implemented via PyLucene 9.4 with Krovetz stemming and default $k_1{=}1.2$, $b{=}0.75$; (2)~\textbf{DPR}~\citep{karpukhin2020dpr}, using the \mbox{\texttt{facebook/dpr-question\_encoder-multiset-base}} checkpoint from HuggingFace, with FAISS~\citep{johnson2021faiss} flat-L2 indexing; (3)~\textbf{ColBERT}~\citep{khattab2020colbert} v2, using the \mbox{\texttt{colbert-ir/colbertv2.0}} checkpoint with PLAID indexing; (4)~\textbf{SPLADE}~\citep{formal2021splade}, using the \mbox{\texttt{naver/splade-cocondenser-ensembledistil}} checkpoint; (5)~\textbf{Hybrid} (BM25+DPR reranked), combining BM25 top-100 with DPR reranking via Reciprocal Rank Fusion~\citep{cormack2009rrf} (RRF $k=60$). All retrievers return top-$k=5$ passages.

\noindent\textit{Generators}: (1)~\textbf{GPT-4o} (\texttt{gpt-4o-2024-08-06}) via the OpenAI API; (2)~\textbf{Claude 3 Haiku} (\texttt{claude-3-haiku-20240307}) via the Anthropic API; (3)~\textbf{Llama-3-8B-Instruct} running locally via \texttt{llama.cpp} with 4-bit GGUF quantization. All generators receive the same prompt template: the question, followed by the top-5 retrieved passages with source URLs, followed by an instruction to answer the question using only the provided passages and to cite each passage used.

\paragraph{Metrics.}
We report: (1) \textbf{Exact Match (EM)}: 1 if the predicted answer string exactly matches the gold span after normalization (lowercase, punctuation removal)~\citep{joshi2017triviaqa}; (2) \textbf{Token F1}: harmonic mean of unigram precision and recall between predicted and gold answer~\citep{kwiatkowski2019nq}; (3) \textbf{Citation Precision} ($P_{\text{cite}}$): fraction of cited passages that contain the gold answer span; (4) \textbf{Citation Recall} ($R_{\text{cite}}$): fraction of questions for which at least one cited passage contains the gold answer span; (5) \textbf{Temporal Degradation Slope} ($\beta_s$): the OLS coefficient from Equation~\ref{eq:degradation}, measuring per-month F1 decay; (6) \textbf{Knowledge Cutoff Gap}: the mean number of days between a question's source document publication date and the evaluation date, serving as a proxy for staleness.

\paragraph{Evaluation protocol.}
We treat months 1--6 as the development window (used for prompt tuning and threshold selection) and months 7--12 as the held-out test window. All reported main results are from the test window. Hyperparameter selection (e.g., the BERTScore grounding threshold $\tau = 0.72$ and the contamination overlap threshold $\theta = 0.4$) was determined on the development window. We report 95\% confidence intervals via bootstrap resampling (1,000 iterations) for all aggregate metrics.

\paragraph{Implementation details.}
All experiments run on a single machine with a 16-core Intel Core i9-13900K CPU and 64 GB RAM; no GPU is required. Dense index construction for DPR uses \texttt{sentence-transformers 2.7.0} and \texttt{faiss-cpu 1.8.0}. ColBERT indexing uses \texttt{colbert-ai/RAGatouille 0.0.8}. BM25 retrieval uses \texttt{rank-bm25 0.2.2}. LLM API calls are cached to disk using the \texttt{diskcache} library; all cached responses are included in the reproducibility package. Total wall-clock time for one monthly evaluation cycle is approximately 4.5 hours. The total API cost for the 12-month evaluation window is \$214 (GPT-4o: \$183, Claude 3 Haiku: \$31).
